\documentclass[conference]{IEEEtran}
\IEEEoverridecommandlockouts

\usepackage{cite}
\usepackage{amsmath,amssymb,amsfonts}
\usepackage{algorithmic}
\usepackage{algorithm}
\usepackage{graphicx}
\usepackage{textcomp}
\usepackage{xcolor}
\usepackage{booktabs}
\usepackage{array}
\usepackage[hyphens]{url}

\def\BibTeX{{\rm B\kern-.05em{\sc i\kern-.025em b}\kern-.08em
    T\kern-.1667em\lower.7ex\hbox{E}\kern-.125emX}}

\newcommand{\sysname}{GhostServe}

\begin{document}

\pdfpagewidth=8.5in
\pdfpageheight=11in

\newcommand{\iscasubmissionnumber}{NaN}

\pagenumbering{arabic}

\title{\sysname{}: Profile-Driven Online Emulation of LLM Inference Serving Systems}
\author{\normalsize{ISCA 2026 MLArchSys Workshop Submission
    \textbf{\#\iscasubmissionnumber} -- Confidential Draft -- Do NOT Distribute!!}}

\maketitle
\thispagestyle{plain}
\pagestyle{plain}

\begin{abstract}
Evaluating scheduling, batching, or capacity decisions on modern LLM
serving engines typically requires the very GPUs those decisions are
meant to conserve. Prior simulators are either offline and
event-driven, hardware-coupled and closed-source, or analytical and
miss the batch-composition dynamics that dominate per-step latency.
We present \sysname{}, an online emulator that runs inside a stock
vLLM process as a runtime plugin: the scheduler, HTTP stack,
admission path, and output pipeline remain real; only the GPU
forward pass is replaced by a profile-driven latency draw. The
emulator is deliberately minimal --- a single density-aware oracle
sampled on each step, and a timer-resolved Future that returns fake
output tokens after the drawn delay --- with no hand-coded mirror of
vLLM internal behaviour. Evaluated across three cells (Qwen2.5-3B on
RTX~8000, Qwen3-8B on RTX~8000, Qwen3-8B on A10) with a
ShareGPT-derived workload at request rates $r\in\{2,4,8,16,32\}$,
\sysname{} tracks the real engine closely on time-to-first-token,
time-per-output-token, inter-token latency, and end-to-end latency
across all sixty (cell, rate, metric) entries. We
discuss the methodology preconditions ---
profile$\leftrightarrow$bench state match and
profile$\leftrightarrow$workload shape match --- that make this
achievable, and the limitations of high-level profile keying at the
single-node regime we evaluate.
\end{abstract}

% ---------------------------------------------------------------------
\section{Introduction}
% ---------------------------------------------------------------------

Continuous-batching LLM serving engines~\cite{kwon2023vllm,yu2022orca}
are the dominant execution substrate for production inference, and
most open research questions on LLM serving---scheduler
design~\cite{agrawal2024sarathi}, prefill-decode
disaggregation~\cite{zhong2024distserve,patel2024splitwise},
replication and placement~\cite{li2023alpaserve,sun2024llumnix},
capacity planning~\cite{patel2024polca}---are phrased as ``what
happens if we change \textit{this} policy?'' Answering such questions
on the real engine consumes the same GPU capacity the policy is
trying to conserve.

Existing alternatives each miss part of the problem. Event-driven
simulators~\cite{agrawal2024vidur,cho2024llmservingsim,
nvidia2026aiconfigurator} re-implement the scheduler and model
forward-pass cost from per-operator latency databases; they lag the
rapidly evolving serving-engine codebase and re-express its policies
in a parallel implementation that must be kept consistent by hand.
Recent emulator work~\cite{agrawal2026revati} addresses the
maintenance burden explicitly via CUDA-API interception, but remains
tied to a fixed hardware profile matrix and requires per-framework
patches. All three approaches ship with a finite supported-device
list (most commonly H100/A100 SXM); the consumer-tier and mid-tier
accelerators (RTX, A10, L40S) typical of research and small-scale
deployment fall outside their matrices.

We occupy a different point in the design space: a
\emph{serving-native} emulator that runs inside a stock vLLM process,
accepts real HTTP traffic, and replaces only the GPU forward pass
with a profile-driven latency draw. The scheduler, tokenizer,
admission path, HTTP layer, output pipeline, and metrics endpoint
are all vLLM's own code, unchanged. Only the forward pass is faked.
The emulator supports two execution modes from the same plugin: a
\emph{realtime} mode that blocks for the drawn latency so online
workloads (live HTTP traffic, dynamic arrival patterns, queueing
under saturation) behave exactly as on real hardware --- our
primary target for \emph{dynamic} serving-system
research~\cite{sun2024llumnix,patel2024polca} --- and an
\emph{accelerated} mode that time-skips (virtual time advances
by the drawn latency without wall-clock blocking, analogous to
Revati~\cite{agrawal2026revati}), enabling offline
configuration-sweep simulation at many-times-real-time speed.

This positioning is conceptually simple but difficult to make
accurate. Two classes of methodology bug consistently invalidate
naive implementations: profile captures that contaminate per-rate
statistics through accumulated server state, and emulator-vs.-real
validations that compare subtly different workloads. A third, more
subtle, trap is the temptation to add compensating ``mirror'' code
that mimics vLLM's internal behaviour on the emulator side---scheduler
admission proxies, CPU-side bookkeeping reimplementations---because
they appear to reduce error against a specific workload. We show that
such mirrors accumulate workload-specific biases that invert when the
workload shape changes, and that a profile-driven design can avoid
them entirely.

\sysname{}'s design is minimal: a density-aware oracle whose pooling
radius grows only where sample density is insufficient for a reliable
estimate, and a timer-resolved Future that returns fake output tokens
after the oracle's predicted delay. Across three cells covering three
model scales and two GPU classes at request rates
$r\in\{2,4,8,16,32\}$, this suffices: every (cell, rate, metric)
entry tracks the real engine closely (\S\ref{sec:results}).

\textbf{Contributions.}
\begin{enumerate}
  \item A plugin-shaped online emulator for vLLM with a small
    integration footprint and no vLLM-source modification
    (\S\ref{sec:design}).
  \item Two reproducibility invariants
    (profile$\leftrightarrow$bench server-state match, including
    identical vLLM CLI arguments; and profile$\leftrightarrow$workload
    shape match) that we argue are preconditions, not tuning knobs,
    for any accuracy claim (\S\ref{sec:design-invariants}).
  \item A density-aware 2D oracle with a statistical sample-count
    reliability floor that breaks a Pareto tradeoff between
    accuracy regimes without introducing a hand-tuned pooling
    constant (\S\ref{sec:design-oracle}).
  \item A multi-hardware, multi-model single-node evaluation
    on Qwen2.5-3B and Qwen3-8B across RTX~8000 and A10 GPUs with a
    ShareGPT workload at rates $r\in\{2,4,8,16,32\}$, including the
    A10 cell at $r=32$ which operates at real-server delivery
    capacity (\S\ref{sec:results}).
\end{enumerate}

% ---------------------------------------------------------------------
\section{Background}
\label{sec:background}
% ---------------------------------------------------------------------

\paragraph{Serving-step anatomy.} A single scheduler step in a
continuous-batching engine admits new requests from a pending queue,
forms a batch (optionally via chunked prefill~\cite{agrawal2024sarathi}),
executes a forward pass on the worker process over an inter-process
boundary, samples output tokens, and delivers them back to the
per-request HTTP streams. Per-step wall-clock time is jointly
determined by batch composition (total tokens in the step) and
concurrency (number of running requests).

\paragraph{Metrics.} We report three per-request metrics: time to
first token (TTFT), time per output token between successive tokens
(TPOT), and end-to-end latency (E2E), plus the inter-token latency
(ITL). TTFT captures admission plus prefill; TPOT and ITL capture
steady-state decode behaviour; E2E accumulates both.

\paragraph{Real-user workload.} We use ShareGPT~\cite{sharegpt} as the
primary evaluation workload. ShareGPT is a public dump of
conversational prompts and responses; the first-turn human prompt of
each conversation serves as the request, and the gpt response length
as the target output length. Conversations with prompts longer than
256 tokens or responses longer than 128 tokens are filtered out to
match the capacity envelope of a single Qwen3-8B deployment on an
RTX~8000. The resulting distribution is heavy-tailed: mean input
length 49 tokens, median 28, p99 227.

% ---------------------------------------------------------------------
\section{Design}
\label{sec:design}
% ---------------------------------------------------------------------

\begin{figure}[t]
  \centering
  \includegraphics[width=\columnwidth]{../../figures/fig1_emulator_architecture.png}
  \caption{\sysname{} integrates at the executor boundary of a stock
  vLLM process. Scheduler, admission path, HTTP server, tokenizer, and
  output pipeline are vLLM's own code, unchanged. A single plugin
  replaces the worker's forward pass with a profile-driven delayed
  Future: the oracle draws a per-step latency from a two-dimensional
  profile keyed by batch composition (total tokens) and concurrency;
  a timer resolves the Future with fake output tokens after the drawn
  delay. No admission proxy, no CPU-side bookkeeping mirror.}
  \label{fig:arch}
\end{figure}

\subsection{Overview}
Figure~\ref{fig:arch} shows the integration. A single environment
variable enables the plugin; when disabled, the source tree runs an
unmodified vLLM. When enabled, the executor's per-step entry point is
redirected to a three-component path: an oracle that returns a
forward-pass latency drawn from a 2D profile, and a
\texttt{threading.Timer} that resolves a Future after that delay,
preserving vLLM's async engine semantics (Figure~\ref{fig:timeline}).
No admission-path proxy, no CPU-side bookkeeping
reimplementation---admission uses vLLM's own queue-and-drain code,
output processing uses vLLM's own sample-then-stream code, and the
oracle's step-cycle sample already includes whatever scheduler-worker
IPC cost real vLLM incurs at that (tt, conc) bucket.

\begin{figure}[t]
  \centering
  \includegraphics[width=\columnwidth]{../../figures/fig2_timer_future_timeline.png}
  \caption{Timer-based Future preserves the engine's scheduler--worker
  overlap. On the real path, the scheduler prepares the next batch
  concurrently with worker execution. The emulated path resolves a
  Future after the profiled delay, so the scheduler continues to
  overlap with the ``pending'' window; downstream pipelining is
  bit-identical.}
  \label{fig:timeline}
\end{figure}

\subsection{Reproducibility invariants}
\label{sec:design-invariants}

Two preconditions are non-negotiable for any accuracy measurement to
be meaningful. Both were observed as bugs in earlier iterations, at
$10$--$15$ percentage points of apparent accuracy each, before the
current methodology stabilised.

\paragraph{I1: profile$\leftrightarrow$bench server-state match.}
Every rate swept during profile capture executes against a freshly
started server, and every vLLM CLI argument used during profile
capture (\texttt{max-num-seqs}, \texttt{max-model-len},
\texttt{enforce-eager}, \texttt{enable-prefix-caching}, \ldots) is
passed identically to the emulated server at validation time. Both
halves of this invariant have been observed as multi-percentage-point
bugs in earlier iterations: a single-session sweep that accumulated
scheduler counters and KV-cache fragmentation caused double-digit
drift across successive days; and a silent mismatch in
\texttt{max-num-seqs} (real configured $=64$ while emu defaulted
$=256$) caused the emulator to run at four times real's allowed
concurrency and diverge catastrophically in the saturation regime,
which was initially misattributed to a missing admission model. A
canonical cleanup-and-warmup harness shared between profile capture
and the validation bench enforces both halves of this invariant and
additionally ensures seed- and parameter-matched workloads on both
sides of the comparison.

\paragraph{I2: profile$\leftrightarrow$workload shape match.} The
profile's sample distribution is specific to the workload shape
captured. A profile captured under one input-length distribution
cannot accurately emulate a workload with a materially different
input-length distribution: long-prompt samples' step cycle differs
from short-prompt samples' step cycle even at the same (tt, conc)
bucket, because batch composition differs (one long prefill vs
several short prefills contributing to the same total-token count).
We therefore profile each target workload separately. This makes
profile capture an up-front per-deployment cost, discussed in
\S\ref{sec:limitations}.

\subsection{Profile pack}
\label{sec:design-profile}
The profile pack is a JSON artifact capturing step-cycle latency as
two joint distributions---prefill and decode---over two-dimensional
buckets keyed by batch composition (total tokens in the step) and
concurrency (number of running requests). Each bucket stores the raw
list of observed latencies rather than a pre-aggregated summary: the
oracle resamples at query time, preserving real variance.

The sweep covers rates from light load to saturation, with denser
sampling in the batching-transit regime where batch composition is
most volatile and oracle queries occur most frequently. Bucket
densities in the transit concurrency band are typically several-fold
higher than in the edges, which is directly responsible for the
accuracy the oracle can achieve at moderate rates.

\subsection{Density-aware oracle}
\label{sec:design-oracle}

Given a query $(t, c)$ at runtime (batch total-tokens, current
concurrency), the oracle returns a latency by sampling from nearby
profile buckets. Two naive choices fail in complementary ways: pure
nearest-neighbor is accurate where buckets are densely populated but
noisy where they are sparse; fixed-$k$ pooling is the reverse---it
smooths over sparse regions but introduces bias from adjacent
outlier samples in dense regions.

\sysname{}'s oracle (Algorithm~\ref{alg:adaptive-k}) makes the
pooling radius depend on local density. Starting from the single
nearest bucket, it expands the neighbor set in 2D-distance order
until cumulative sample count meets a statistical reliability floor
$M$, then returns a Shepard inverse-distance-weighted draw. When the
nearest bucket already contains $\ge M$ samples---the common case in
dense regions---expansion does not happen and behaviour is identical
to pure nearest-neighbor. Only sparse regions expand. We use
$M\!=\!30$, the conventional central-limit threshold for
sample-mean reliability of a fixed-variance estimator; it is not
tuned per hardware or model.

\begin{algorithm}[t]
\caption{Density-aware neighbor pooling}
\label{alg:adaptive-k}
\begin{algorithmic}[1]
\STATE \textbf{input:} query $(t,c)$, bucket set $B$, reliability floor $M$
\STATE sort $B$ by range-normalized 2D distance to $(t,c)$
\STATE $\mathcal{S} \leftarrow \emptyset$,\;\; $n \leftarrow 0$
\FOR{$B_i$ in sorted order}
  \STATE $\mathcal{S} \leftarrow \mathcal{S} \cup \{B_i\}$
  \STATE $n \leftarrow n + |B_i.\text{samples}|$
  \IF{$n \ge M$} \STATE \textbf{break} \ENDIF
\ENDFOR
\RETURN Shepard-weighted sample over $\mathcal{S}$
\end{algorithmic}
\end{algorithm}

% ---------------------------------------------------------------------
\section{Evaluation}
\label{sec:results}
% ---------------------------------------------------------------------

\paragraph{Setup.} Three cells, all vLLM 0.18.1, FP16, maximum
sequence length 4096, chunked-prefill enabled (vLLM default), CUDA
graph enabled (vLLM default): (i)~\textbf{M1b}: Qwen2.5-3B on a
single NVIDIA Quadro RTX~8000 (48\,GB); (ii)~\textbf{M2}: Qwen3-8B
on the same RTX~8000; (iii)~\textbf{M4}: Qwen3-8B on a single
NVIDIA A10 (24\,GB) with \texttt{max-num-seqs=64} to fit the smaller
memory budget. Benchmark: Poisson-arrival workload over
ShareGPT-derived prompts at request rates $r\in\{2,4,8,16,32\}$,
$2000$ prompts per rate, fixed random seed. Each rate is run against
a freshly started server with a shared warmup. The comparison is
between a real run and an emulated run of the identical workload
under the identical harness.

\subsection{Accuracy}
\label{sec:results-main}

Table~\ref{tab:main} reports the per-cell relative error
$(\text{emu}-\text{real})/\text{real}$ for each metric at each rate,
across all three cells and five rates. The A10 cell at $r=32$
operates at the server's sustained delivery capacity ($\approx
16.5$\,r/s, well below the requested $32$\,r/s) because the smaller
A10 is configured with \texttt{max-num-seqs=64}. With that
configuration propagated identically to the emulator's vLLM (per
invariant I1), the density-aware oracle reproduces real-server
behaviour in the saturation regime without any additional admission
mechanism: emu and real both complete at the same rate because both
are concurrency-limited by the same
\texttt{max-num-seqs}.

\begin{table}[t]
\centering
\caption{Relative error $(\text{emu}-\text{real})/\text{real}$
across four cells at $r\in\{2,4,8,16,32\}$. M1b/M2/M4 are the main
accuracy claims (default configuration); R3 is an ablation with
vLLM's prefix-cache disabled. All use the same pipeline:
density-aware adaptive profile capture with saturation-weighted
rate coverage, dense profile, hook OFF, surrogate OFF.}
\label{tab:main}
\small
\begin{tabular}{@{}lrrrrr@{}}
\toprule
Metric & $r{=}2$ & $r{=}4$ & $r{=}8$ & $r{=}16$ & $r{=}32$ \\
\midrule
\multicolumn{6}{@{}l}{\textbf{M1b: Qwen2.5-3B / RTX~8000 / ShareGPT}} \\
TTFT & $-6.43\%$ & $-5.37\%$ & $-6.50\%$ & $-2.48\%$ & $+1.64\%$ \\
TPOT & $+0.03\%$ & $-0.05\%$ & $+0.07\%$ & $-0.18\%$ & $+1.41\%$ \\
ITL  & $-0.09\%$ & $-0.21\%$ & $+0.20\%$ & $-0.41\%$ & $+1.64\%$ \\
E2E  & $+2.67\%$ & $+3.73\%$ & $+2.81\%$ & $+4.25\%$ & $+4.78\%$ \\
\midrule
\multicolumn{6}{@{}l}{\textbf{M2: Qwen3-8B / RTX~8000 / ShareGPT}} \\
TTFT & $-7.88\%$ & $-5.75\%$ & $-1.12\%$ & $+2.13\%$ & $-3.54\%$ \\
TPOT & $+0.17\%$ & $+0.17\%$ & $+0.42\%$ & $+0.67\%$ & $+1.67\%$ \\
ITL  & $-0.09\%$ & $+0.15\%$ & $-0.23\%$ & $+0.49\%$ & $+1.25\%$ \\
E2E  & $-0.84\%$ & $-0.03\%$ & $+0.01\%$ & $-0.32\%$ & $+1.72\%$ \\
\midrule
\multicolumn{6}{@{}l}{\textbf{M4: Qwen3-8B / A10 / ShareGPT}} \\
TTFT & $-2.50\%$ & $-0.56\%$ & $-2.52\%$ & $-2.35\%$ & $-0.45\%$ \\
TPOT & $+1.60\%$ & $+1.90\%$ & $+1.31\%$ & $+1.09\%$ & $+0.28\%$ \\
ITL  & $+0.96\%$ & $+0.89\%$ & $+0.95\%$ & $-0.38\%$ & $-0.74\%$ \\
E2E  & $+1.44\%$ & $+1.57\%$ & $+1.61\%$ & $+0.29\%$ & $-0.95\%$ \\
\midrule
\multicolumn{6}{@{}l}{\textbf{R3: Qwen3-8B / RTX~8000 / ShareGPT, prefix-cache OFF}} \\
TTFT & $-9.01\%$ & $-7.37\%$ & $-4.29\%$ & $-2.83\%$ & $-3.76\%$ \\
TPOT & $+0.05\%$ & $+0.16\%$ & $+0.13\%$ & $+0.81\%$ & $+1.13\%$ \\
ITL  & $-0.14\%$ & $+0.11\%$ & $-0.06\%$ & $+0.46\%$ & $+0.63\%$ \\
E2E  & $-0.43\%$ & $-0.66\%$ & $-0.43\%$ & $-0.72\%$ & $+3.73\%$ \\
\bottomrule
\end{tabular}
\end{table}

\paragraph{Summary.} The three main-configuration cells
(M1b, M2, M4) each pass the $|\Delta|\le 6\%$ gate on the
steady-state metrics (TPOT, ITL, E2E) at every rate; TTFT is
within $|\Delta|\le 10\%$ at every rate. The R3 ablation
(prefix-cache OFF) preserves the same accuracy envelope at
every rate, showing the pipeline is not specific to the default
vLLM configuration. TPOT, ITL, and E2E are the tighter group;
TTFT carries more variance because first-token latency mixes
admission-path timing (which depends on the real scheduler's
iteration cadence) with the first prefill step --- the oracle's
step-cycle sample captures the latter, while vLLM's own
admission code contributes the former without any
emulator-side proxy. Validation here is on the filtered
ShareGPT workload (prompt + output capped at $256+128$~tokens);
extension to long-context workloads is discussed in
\S\ref{sec:limitations} as future work.

% ---------------------------------------------------------------------
\section{Limitations and future work}
\label{sec:limitations}
% ---------------------------------------------------------------------

\paragraph{Profile capture is per-deployment and short-context.}
Each (hardware, model, workload) bundle requires a dedicated
capture. We evaluate on a filtered-length ShareGPT workload
(${\leq}256$~prompt / ${\leq}128$~output tokens) to keep the
profile-capture development cycle short: a dense capture at
$r\in\{2,4,8,16,32\}$ takes ${\approx}2$~hours of GPU time,
versus days for per-operator databases (Vidur, AIConfigurator,
Revati). Extending to long-context workloads (full ShareGPT,
vLLM's default benchmark configuration) is natural future work:
at long prompt lengths, the $(tt, \text{conc})$ bucket grid
becomes sparse at high concurrency and a longer capture
schedule is needed to fill the saturation regime without
aliasing to overloaded-server artifacts. Finer oracle axes
(per-request prompt length rather than batch-total tokens)
could further reduce profile cost for variable-shape workloads.

\paragraph{Single-node scope.} We validate single-GPU
single-process deployments. Multi-GPU parallelism (TP, PP, EP)
and cluster-level placement introduce dynamics the current
oracle does not capture: collective communication latency, NIC
bandwidth sharing, inter-host queue coupling. Extending the
oracle with inter-node round-trip and collective-bandwidth
terms, profiled from a real multi-host measurement, is the
primary direction for capacity-planning workflows.

\paragraph{High-level keying elides low-level variance.} The
(total-tokens, concurrency) oracle does not expose prefix-cache
hit/miss effects, KV-cache offloading under memory pressure,
kernel-launch jitter near peak utilisation, or
speculative-decoding acceptance variance. Per-phenomenon profile
axes are feasible but multiply capture cost super-linearly;
selective additions tied to deployment scenarios are the
practical path.

\paragraph{Integration bindings track vLLM internals.} The
footprint is small ($\sim$400~lines of plugin code, no
source-level modification of vLLM, no reimplementation of
admission, output processing, or worker-side bookkeeping), but
the executor-boundary hook and platform plugin are pinned to
specific vLLM internal APIs; results here are against vLLM
0.18.1. Unlike a scheduler-reimplementation simulator, this is a
small surface of integration points that must stay in sync with
upstream; an automated test matrix across vLLM releases would
catch API drift early.

% ---------------------------------------------------------------------
\section{Related work}
\label{sec:related}
% ---------------------------------------------------------------------

\textbf{LLM serving simulators and emulators.}
Vidur~\cite{agrawal2024vidur} is an offline Python event-driven
simulator that re-implements vLLM's scheduler as \texttt{sarathi}
and predicts per-operator latencies via random-forest models
trained on operator-level microbenchmarks. Its published
hardware matrix is limited to datacenter SXM configurations
(A100 DGX, H100 DGX, 4$\times$A100 NVLink, 8$\times$A40);
consumer-class (Turing, RTX) and mid-tier devices (A10, L40S)
are not supported without a fresh profiling campaign.
LLM-ServingSim~\cite{cho2024llmservingsim} targets disaggregated
serving with an sklearn-based attention-latency predictor and
supports A6000, H100, and TPU-v6e-1.
AIConfigurator~\cite{nvidia2026aiconfigurator} is an NVIDIA
open-source configuration recommender built on an operator-level
profile database and a piecewise-linear saturation correction
factor; its supported device list is H100/H200/B200/GB200/A100
SXM only. Revati~\cite{agrawal2026revati}, from the same group as
Vidur, acknowledges that ``existing discrete-event simulation
approaches face a perpetual maintenance burden'' from
re-implementing serving logic and proposes CUDA-API interception
with a virtual-time coordination protocol; the paper reports
results on vLLM and SGLang but does not publicly release code.

\sysname{} differs on six axes:
(i) \emph{runtime integration} --- runs as a vLLM plugin with zero
framework modifications, leveraging the engine's own scheduler,
admission, and output-processing paths unchanged;
(ii) \emph{dual execution modes from one plugin} --- the same
profile-driven oracle supports both a \emph{realtime} mode for
online/dynamic serving study (live HTTP, genuine queueing, SLA
measurement under load) and an \emph{accelerated} time-skip mode
for offline configuration-sweep simulation (as in Vidur and
Revati). Vidur, AIConfigurator, and LLM-ServingSim offer only the
offline path; Revati is emulator-only. \sysname{} covers both
from a single implementation;
(iii) \emph{maintenance burden} --- no re-implementation of
serving logic; the vLLM v0$\to$v1 architectural transition is
absorbed automatically without code changes on our side, whereas
Vidur's \texttt{sarathi} engine and AIConfigurator's operator
database require manual porting, and Revati's LD\_PRELOAD patches
must be re-synced per vLLM release;
(iv) \emph{profile granularity} --- a single aggregate
step-cycle capture (${\approx}2$~hours per hardware/model pair)
rather than per-operator profiling (days);
(v) \emph{backend- and hardware-portable by construction} ---
the profile captures whatever attention kernel and KV-cache
backend vLLM executes on whatever GPU is present, giving
validated coverage across Turing (RTX~8000) and Ampere (A10), two
devices absent from the above simulators' supported matrices;
(vi) \emph{no learned hyperparameters} --- the oracle is empirical
distribution lookup with adaptive $K$ and a CLT-based sample
floor, not random-forest training. All profile-based approaches
(including ours) require the target hardware to be physically
available for profile capture. We plan to release \sysname{} as
an open-source pip-installable plugin concurrent with this
publication.

\textbf{Serving systems and scheduling.} \sysname{} plugs into
vLLM~\cite{kwon2023vllm}. Adjacent systems whose policies would
benefit from accurate emulation include
Sarathi-Serve~\cite{agrawal2024sarathi},
DistServe~\cite{zhong2024distserve} and
Splitwise~\cite{patel2024splitwise} for prefill/decode
disaggregation, AlpaServe~\cite{li2023alpaserve} for model-parallel
placement, Llumnix~\cite{sun2024llumnix} for dynamic rescheduling,
and POLCA~\cite{patel2024polca} for capacity planning.

\textbf{Continuous batching.} Orca~\cite{yu2022orca} introduced
iteration-level scheduling, which vLLM~\cite{kwon2023vllm} scaled to
production via paged attention.
FlashAttention~\cite{dao2022flashattention} and
FlashInfer~\cite{ye2025flashinfer} provide the attention kernels
\sysname{}'s profile implicitly characterises.

% ---------------------------------------------------------------------
\section{Conclusion}
\label{sec:conclusion}
% ---------------------------------------------------------------------

\sysname{} is an online, serving-native LLM inference emulator that
runs as a plugin inside a stock vLLM process with no source
modification. Its design deliberately keeps the emulator's surface
area small: a single density-aware oracle draws from a 2D
step-cycle profile, a timer-resolved Future preserves vLLM's async
semantics, and vLLM's own code handles admission, batching, output
processing, and the HTTP path. Across three cells spanning two GPU
classes (RTX~8000, A10) and two model scales (Qwen2.5-3B, Qwen3-8B)
with a ShareGPT workload at $r\in\{2,4,8,16,32\}$, this suffices to
track the real engine closely across all 60 (cell, rate, metric)
entries, including the saturation-limited A10 cell at $r=32$. The
accuracy rests on two reproducibility invariants we argue are
preconditions rather than tuning knobs. The limitations are scope:
single-node evaluation, workload-specific profiles, and high-level
keying that hides some low-level variance---each a concrete
future-work direction rather than a methodological dead end.

%%%%%%%%% -- BIB STYLE AND FILE -- %%%%%%%%
\bibliographystyle{IEEEtranS}
\bibliography{refs}

\end{document}
